\documentclass[8pt, a4paper]{article}

% ── Packages ─────────────────────────────────────────────────────────────────
\usepackage[a4paper, top=2.5cm, bottom=2.5cm, left=2.8cm, right=2.8cm]{geometry}
\usepackage[T1]{fontenc}
\usepackage[utf8]{inputenc}
\usepackage{lmodern}
\usepackage{microtype}
\usepackage{parskip}          % space between paragraphs, no indent
\usepackage{titlesec}         % section heading styles
\usepackage{enumitem}         % list customisation
\usepackage{booktabs}         % professional tables
\usepackage{tabularx}         % auto-width table columns
\usepackage{array}            % extra column formatting
\usepackage{xcolor}           % colours
\usepackage{mdframed}         % framed boxes
\usepackage{fancyhdr}         % header / footer
\usepackage{hyperref}         % clickable links (load last)

% ── Colours ───────────────────────────────────────────────────────────────────
\definecolor{accentblue}{RGB}{30, 90, 160}
\definecolor{rulegray}{RGB}{190, 195, 205}
\definecolor{boxbg}{RGB}{237, 242, 250}

% ── Hyperref setup ────────────────────────────────────────────────────────────
\hypersetup{
  colorlinks = true,
  linkcolor  = accentblue,
  urlcolor   = accentblue,
}

% ── Section heading styles ────────────────────────────────────────────────────
\titleformat{\section}
  {\large\bfseries\color{accentblue}}
  {\thesection.}{0.6em}{}
  [\vspace{-4pt}\textcolor{rulegray}{\rule{\linewidth}{0.6pt}}\vspace{2pt}]

\titlespacing{\section}{0pt}{18pt}{6pt}

% ── Concept box ───────────────────────────────────────────────────────────────
\newmdenv[
  backgroundcolor   = boxbg,
  linecolor         = accentblue,
  linewidth         = 1.2pt,
  leftline          = true,
  rightline         = false,
  topline           = false,
  bottomline        = false,
  innerleftmargin   = 10pt,
  innertopmargin    = 6pt,
  innerbottommargin = 6pt,
  skipabove         = 6pt,
  skipbelow         = 6pt,
]{conceptbox}

% ── List styling ──────────────────────────────────────────────────────────────
\setlist[itemize]{topsep=3pt, itemsep=2pt, parsep=0pt, leftmargin=1.6em}
\setlist[enumerate]{topsep=3pt, itemsep=2pt, parsep=0pt, leftmargin=1.8em}

% ── Header / Footer ───────────────────────────────────────────────────────────
\pagestyle{fancy}
\fancyhf{}
\renewcommand{\headrulewidth}{0.4pt}
\fancyhead[L]{\small\color{accentblue}\textbf{ELEC3609 / ELEC9609}}
\fancyhead[R]{\small\color{gray}Week 2 --- Web Fundamentals \& Project Management}
\fancyfoot[C]{\small\thepage}

% =============================================================================
\begin{document}

% ── Title Block ───────────────────────────────────────────────────────────────
\begin{center}
  {\LARGE\bfseries\color{accentblue} ELEC3609 / ELEC9609: Internet Software Platforms}\\[6pt]
  {\large Week 2 Study Guide: Web Fundamentals \& Project Management}
\end{center}

\vspace{12pt}
\textcolor{rulegray}{\rule{\linewidth}{1pt}}
\vspace{6pt}

% =============================================================================
\section{The End-to-End Web Request Lifecycle}

When a user types a URL (e.g., \texttt{https://cusp.sydney.edu.au/students}) into a browser and hits Enter, a multi-step process is triggered across the client, network, and server stack.

\begin{conceptbox}
\textbf{Client Side (Browser \& OS)}
\begin{enumerate}
  \item \textbf{URL Parsing:} Browser breaks URL into scheme, domain, port, and path.
  \item \textbf{Cache Check:} Browser checks local cache for pre-saved static assets.
  \item \textbf{Network Check:} OS identifies local network configuration via DHCP.
  \item \textbf{DNS Resolution:} OS resolves host domain name into an IP address.
  \item \textbf{TCP Connection:} OS establishes a TCP connection to the server (Port 80/443).
  \item \textbf{HTTP Request:} Browser constructs and sends an HTTP/HTTPS request.
\end{enumerate}
\end{conceptbox}

\begin{conceptbox}
\textbf{Server Side (Infrastructure \& Back-End)}
\begin{enumerate}
  \setcounter{enumi}{6}
  \item \textbf{Load Balancing:} Load Balancer / WAF receives TCP traffic, terminates SSL/TLS, and routes HTTP request to an available App Server.
  \item \textbf{Business Logic:} App Server processes request logic and queries Database(s).
  \item \textbf{Response Render:} App Server inserts DB data into templates and generates HTTP response.
\end{enumerate}
\end{conceptbox}

\begin{conceptbox}
\textbf{Client Rendering (Browser)}
\begin{enumerate}
  \setcounter{enumi}{9}
  \item \textbf{Parse Response:} Browser reads status code (e.g., 200 OK) and headers.
  \item \textbf{Asset Fetching:} Browser parses HTML into DOM tree; fetches external CSS, JS, images.
  \item \textbf{Screen Render:} Rendering Engine builds final UI for user display.
\end{enumerate}
\end{conceptbox}

\medskip

\begin{conceptbox}
\textbf{URL Anatomy Breakdown}

\smallskip
Example URL: \texttt{https://cusp.sydney.edu.au/students/view-unit-page/uos\_id/101257}

\smallskip
\begin{tabularx}{\linewidth}{l l l X}
  \toprule
  \textbf{Component} & \textbf{Example} & \textbf{Technical Role} & \textbf{Description} \\
  \midrule
  Protocol  & \texttt{https://} & Scheme              & HyperText Transfer Protocol Secure \\
  Subdomain & \texttt{cusp}     & Subdomain           & Host division within main domain \\
  Domain    & \texttt{sydney}   & Organization Domain & Second-level domain identifier \\
  gTLD      & \texttt{.edu}     & Generic TLD         & Top-level category (Education) \\
  ccTLD     & \texttt{.au}      & Country Code TLD    & Top-level country domain (Australia) \\
  Path      & \texttt{/stud...} & Resource Path       & Server file/route path \\
  \bottomrule
\end{tabularx}
\end{conceptbox}

% =============================================================================
\section{Network Configuration, DNS \& OSI Transport Layer}

\begin{conceptbox}
\textbf{DHCP (Dynamic Host Configuration Protocol)}

When joining a network, DHCP configures host network parameters:
\begin{enumerate}
  \item \textbf{Local IP Address:} Unique identity of computer on local network.
  \item \textbf{Subnet Mask:} Defines local network boundary vs.\ external destinations.
  \item \textbf{Default Gateway:} Local router interface that forwards non-local network traffic.
  \item \textbf{Nameservers (DNS):} IP addresses of DNS servers used to resolve domain names.
\end{enumerate}
\end{conceptbox}

\begin{conceptbox}
\textbf{DNS (Domain Name System) Resolution Flow}

Connections require numeric IP addresses, not domain strings. Resolution sequence:
\begin{enumerate}
  \item Check OS Local DNS Entries (e.g., \texttt{/etc/hosts}).
  \item Identify Nameservers (e.g., \texttt{/etc/resolv.conf} configured via DHCP).
  \item Submit UDP request on Port 53 to DNS nameserver.
  \item \textbf{Recursive Lookup:} Nameservers recursively ask parent DNS servers until the authoritative server returns the IP matching the domain's A-record.
\end{enumerate}
\end{conceptbox}

\begin{conceptbox}
\textbf{OSI Model Layers (Key Reference Points)}
\begin{itemize}
  \item \textbf{Layer 1 (Physical):} Copper cables, Fiber optics, Coaxial.
  \item \textbf{Layer 2 (Data Link):} Ethernet, Wi-Fi (IEEE 802.11a/b/g/n).
  \item \textbf{Layer 3 (Network):} IP protocol, routing between hosts.
  \item \textbf{Layer 4 (Transport):} TCP / UDP data transfer protocols.
  \item \textbf{Layer 7 (Application):} HTTPS, HTTP, SSH, DNS services.
\end{itemize}
\end{conceptbox}

\begin{conceptbox}
\textbf{Transport Layer Protocols: TCP vs.\ UDP}

\smallskip
\begin{tabularx}{\linewidth}{l X X}
  \toprule
  \textbf{Feature} & \textbf{TCP (Transmission Control Protocol)} & \textbf{UDP (User Datagram Protocol)} \\
  \midrule
  Data Unit          & Segments                                      & Datagrams \\
  Connection         & Connection-oriented (3-way handshake)         & Connectionless (No handshake) \\
  Delivery Guarantee & Guaranteed, in-order delivery                 & ``Best effort'', no delivery context \\
  Error Recovery     & Retransmits missing/corrupted data            & No retransmission or ACK \\
  Default Ports      & Port 80 (HTTP), Port 443 (HTTPS)              & Port 53 (DNS) \\
  Attack Vector      & TCP SYN Flooding                              & UDP Spoofing / Amplification \\
  \bottomrule
\end{tabularx}
\end{conceptbox}

% =============================================================================
\section{Server-Side Architecture \& Database Design}

\begin{conceptbox}
\textbf{Typical 3-Tier Web Architecture}

\medskip
\textbf{Tier 1: Internet-Facing / Front-End Edge Services}
\begin{itemize}
  \item WAF / DDoS Protection (e.g., Cloudflare)
  \item Content Delivery Network (CDN) for static assets
  \item Elastic Load Balancer (ELB): SSL Offloading \& Routing
\end{itemize}

\textbf{Tier 2: Application Servers (Back-End Logic)}
\begin{itemize}
  \item Workflows, APIs, Business Logic, Auth
  \item Frameworks: Django (Python), Express (JS), Rails
  \item Hosting: EC2 (VMs), Docker (Containers), Lambda
\end{itemize}

\textbf{Tier 3: Persistent Database Layer (Relational / NoSQL)}
\begin{itemize}
  \item Master Database (Primary Writes) $\longrightarrow$ Slave Replicas (Read-Only Scale)
  \item App Servers write to Master; reads distributed across Slaves.
\end{itemize}
\end{conceptbox}

\begin{conceptbox}
\textbf{Database Systems: SQL vs.\ NoSQL}

\medskip
\textbf{1.\ SQL (Relational Databases):}
\begin{itemize}
  \item \textbf{Structure:} Strict relational tables, predefined schemas, explicit primary/foreign keys.
  \item \textbf{Strengths:} Data integrity, complex JOIN queries, ACID transaction compliance.
  \item \textbf{Examples:} MySQL, PostgreSQL, SQLite, MS SQL Server.
\end{itemize}

\textbf{2.\ NoSQL (Non-Relational Databases):}
\begin{itemize}
  \item \textbf{Structure:} Flexible schemas, key-value stores, JSON documents, graph objects.
  \item \textbf{Strengths:} High throughput, scalable horizontal partitioning, dynamic data structures.
  \item \textbf{Examples:} MongoDB (Document), Redis (Key-Value/Cache), DynamoDB (Cloud Key-Value).
\end{itemize}
\end{conceptbox}

\begin{conceptbox}
\textbf{Database Replication: Master/Slave \& Slave Drift}
\begin{itemize}
  \item \textbf{Master (Primary):} Accepts both READ and WRITE operations. Single primary master prevents write conflicts.
  \item \textbf{Slaves (Replicas):} Accept READ operations ONLY. Replicates master changes asynchronously for read-scaling.
\end{itemize}

\medskip
\textbf{EXAM CRITICAL: Slave Drift (Replication Delay)}
\begin{itemize}
  \item \textbf{Definition:} The temporary time delay between a write execution on the Master DB and propagation to Slave DBs.
  \item \textbf{Vulnerability:} Reading from a Slave DB immediately after writing to the Master DB can return stale data.
  \item \textbf{Risk:} Can cause double billing, inventory over-selling, or authorization race conditions.
  \item \textbf{Solution:} High-stakes business decisions (payments, auth, critical updates) MUST read directly from the Master DB!
\end{itemize}
\end{conceptbox}

\begin{conceptbox}
\textbf{Web Server Software \& Technology Stacks}

\medskip
\textbf{Server Engines:}
\begin{itemize}
  \item \textbf{Nginx:} Event-driven worker processes; lightweight, low RAM footprint, highly scalable reverse proxy.
  \item \textbf{Apache:} Process/thread-per-request architecture; easy setup but higher memory overhead under heavy load.
  \item \textbf{IIS:} Microsoft web server; historically higher attack vulnerability, less common for open-source web apps.
\end{itemize}

\textbf{Common Software Stacks:}
\begin{itemize}
  \item \textbf{LAMP:} Linux (OS) + Apache (Server) + MySQL (DB) + PHP (Language)
  \item \textbf{MEAN:} MongoDB (DB) + Express.js (Back-end) + Angular.js (Front-end) + Node.js (Runtime)
  \item \textbf{LYME:} Linux + Yaws + Mnesia/CouchDB + Erlang
\end{itemize}
\end{conceptbox}

% =============================================================================
\section{HTTP Protocol, Status Codes \& Client-Side Rendering}

\begin{conceptbox}
\textbf{HTTP Message Structure}
\begin{itemize}
  \item \textbf{Header:} Contains metadata (Cookies, Content-Type, Content-Length, Auth Tokens, Status Codes).
  \item \textbf{Body:} Contains actual payload (HTML document, JSON API payload, images). Optional in GET/HEAD requests.
\end{itemize}
\end{conceptbox}

\begin{conceptbox}
\textbf{HTTP Request Methods}
\begin{itemize}
  \item \textbf{GET:} Requests resource data from server (Headers only, no request body).
  \item \textbf{POST:} Sends data payload to server (Forms, file uploads, dynamic inputs).
  \item \textbf{HEAD:} Same as GET, but server returns ONLY response headers (omits response body).
  \item \textbf{PUT:} Replaces/creates a target resource completely at destination URL.
  \item \textbf{DELETE:} Removes specified target resource from server.
  \item \textbf{OPTIONS:} Queries server to discover permitted HTTP request methods on a resource.
\end{itemize}
\end{conceptbox}

\begin{conceptbox}
\textbf{HTTP Response Status Codes}

\smallskip
\begin{tabularx}{\linewidth}{l l X}
  \toprule
  \textbf{Code} & \textbf{Category} & \textbf{Common Examples \& Descriptions} \\
  \midrule
  1xx & Informational & 100 Continue, 101 Switching Protocols \\
  2xx & Success       & 200 OK (Standard success response) \\
  3xx & Redirection   & 301 Moved Permanently, 302 Found (Temp Redirect), 304 Not Modified (Cached copy valid) \\
  4xx & Client Error  & 400 Bad Request, 403 Forbidden, 404 Not Found \\
  5xx & Server Error  & 500 Internal Server Error, 502 Bad Gateway \\
  \bottomrule
\end{tabularx}
\end{conceptbox}

\begin{conceptbox}
\textbf{DOM Construction \& Browser Rendering}
\begin{itemize}
  \item \textbf{DOM (Document Object Model):} In-memory tree representation of HTML structure constructed by browser.
  \item \textbf{Asset Fetching:} Browser executes additional HTTP requests for external CSS (\texttt{<link>}) and JS (\texttt{<script>}).
  \item \textbf{Browser Rendering Engines:}
  \begin{itemize}
    \item \textbf{Blink:} Google Chrome, Brave, modern MS Edge
    \item \textbf{WebKit:} Apple Safari
    \item \textbf{Gecko:} Mozilla Firefox
    \item \textbf{EdgeHTML:} Legacy MS Edge
  \end{itemize}
\end{itemize}
\end{conceptbox}

% =============================================================================
\section{Project Management \& Requirements Engineering}

\begin{conceptbox}
\textbf{Project Management Tooling}
\begin{itemize}
  \item \textbf{Written Documentation:} Google Docs (Real-time collaboration), \LaTeX{} (Git-friendly, professional formatting), MS Word.
  \item \textbf{Diagramming Tools:} Diagrams.net (Draw.io), Lucidchart, Balsamiq (Wireframing), MS Visio.
  \item \textbf{Workflow Management:} GitHub (Issues/PRs), Jira (Agile/SDLC integration), Phabricator, Asana/Trello, MS Project (Gantt charts).
\end{itemize}
\end{conceptbox}

\begin{conceptbox}
\textbf{Requirements Engineering Process}

\medskip
\textbf{Step 1: Identify Problems} $\to$ Generate explicit problem statements.\\
\textit{Example: ``Sellers have difficulty determining optimal item listing prices.''}

\medskip
\textbf{Step 2: Determine Use Cases} $\to$ Map problem statements into actor actions.
\begin{itemize}
  \item \textbf{Actor:} System role (e.g., Buyer, Admin), NOT an individual person.
  \item \textbf{Action:} Discrete, well-defined task executable by an actor.
\end{itemize}
\end{conceptbox}

\begin{conceptbox}
\textbf{Use Case Document Key Fields}

When documenting use cases in system specifications:
\begin{itemize}
  \item \textbf{Name \& Goal:} Short title and objective.
  \item \textbf{Actors:} Initiating and secondary system roles.
  \item \textbf{Preconditions:} Mandatory conditions prior to use case execution.
  \item \textbf{Main Flow:} Step-by-step primary success scenario.
  \item \textbf{Alternative Flows:} Exception handling and error branches.
  \item \textbf{Priority:} Development ranking for Minimum Viable Product (MVP).
\end{itemize}
\end{conceptbox}

\begin{conceptbox}
\textbf{System Modeling Diagrams}
\begin{itemize}
  \item \textbf{Sequence Diagrams:} Visualizes chronological order of interactions between actors and objects over time.
  \item \textbf{State Diagrams:} Models entity lifetimes by specifying valid states and allowed transitions (e.g., Draft $\to$ Submitted $\to$ Approved).
  \item \textbf{Wireframes:} Low-fidelity visual layouts (``cartoons'') of UI pages. Storyboards link wireframes interactively.
\end{itemize}
\end{conceptbox}

\begin{conceptbox}
\textbf{Software Requirements Specification (SRS)}
\begin{itemize}
  \item \textbf{Definition:} Master agreement specifying system functions, constraints, and non-goals between client and team.
  \item \textbf{Critical Benefit:} Eliminates ``Scope Creep'' by establishing a clear definition of project completion.
\end{itemize}

\medskip
\textbf{SRS Document Structure:}
\begin{enumerate}
  \item \textbf{Introduction:} Purpose, Definitions, System Overview, References.
  \item \textbf{Overall Description:} Product Perspective, User Interfaces, Constraints, Assumptions.
  \item \textbf{Specific Requirements:} Testable granular rules (Functional, External Interfaces, Security, Reliability, Portability).
\end{enumerate}
\end{conceptbox}

% =============================================================================
\section{Exam \& Revision Checklist}

\begin{itemize}
  \item[$\square$] Can you trace a web request from browser URL bar to screen rendering?
  \item[$\square$] What 4 items are assigned to a machine via DHCP?
  \item[$\square$] What is the core functional difference between TCP and UDP?
  \item[$\square$] Why does ``Slave Drift'' occur in DB replication, and how do you protect payment reads?
  \item[$\square$] Compare SQL vs.\ NoSQL with 2 real-world technology examples each.
  \item[$\square$] Identify the HTTP status code range for Client Errors vs.\ Server Errors.
  \item[$\square$] What purpose does an SRS document serve in managing project scope?
  \item[$\square$] What is the difference between a Sequence Diagram and a State Diagram?
\end{itemize}

\vspace{14pt}
\textcolor{rulegray}{\rule{\linewidth}{1pt}}

\end{document}
